% Experimental Setup -- ICRA section, IEEEtran two-column.
%
% Drop into a paper that already loads: amsmath, amssymb, booktabs, siunitx,
% url. Uses \texttt for identifiers and no custom macros, so it can be pasted
% without carrying a preamble.
%
% A standalone article-class wrapper is in experimental-setup-standalone.tex.

\section{Experimental Setup}
\label{sec:setup}

This section specifies the policy, its training, and the mapping and evaluation
pipeline used for the node-placement study. Quantities are given in the units
the implementation uses; where a constant is load-bearing, we state why rather
than only what.

\subsection{Simulation and Task}
\label{sec:setup:task}

Experiments are conducted in CARLA~0.9.15 at Epic rendering quality. The
perception front-end is a frozen autoencoder trained on Epic-quality frames, so
rendering quality is a fixed property of the setup rather than a tunable:
lowering it would shift the input distribution out from under the encoder.

The agent controls a single ego vehicle in Town02. A task instance is a
point-goal: the agent is given a goal pose at a controlled geodesic distance
from its start and must arrive within an arrival radius. Goal distance is drawn
uniformly from a band; training pins this band to \SI{5}{\meter}--\SI{20}{\meter},
and evaluation additionally reports \SI{20}{\meter}--\SI{40}{\meter} and
\SI{40}{\meter}--\SI{70}{\meter}, which are extrapolation with respect to the
training distribution.

The action space is two-dimensional and continuous,
$a = (\mathrm{steer}, \mathrm{throttle})$, applied as absolute commands rather
than increments, with $|\mathrm{steer}| \le 0.5$,
$\mathrm{throttle} \in [0, 0.6]$, and a target speed of \SI{20}{\kilo\meter\per\hour}.
The observation records the control \emph{actually applied} to the simulator
rather than the raw action, since under incremental control the two differ.

\subsection{Observation Space}
\label{sec:setup:obs}

The policy observes a dictionary of four keys totalling 145 dimensions
(Table~\ref{tab:obs}). The keys are kept separate rather than concatenated into
one flat vector so that each is projected independently; flattening
heterogeneous quantities behind a single normalisation layer mixes incomparable
scales.

\begin{table}[t]
\centering
\caption{Observation space, \num{145} dimensions total.}
\label{tab:obs}
\begin{tabular}{@{}llp{4.6cm}@{}}
\toprule
Key & Dim. & Contents \\
\midrule
\texttt{latent} & 128 & frozen autoencoder embedding of the
                        $120\!\times\!120$ RGB frame \\
\texttt{goal}   & 12  & egocentric goal offset (2), heading
                        $(\sin\psi, \cos\psi)$ (2), action history (8) \\
\texttt{vector} & 4   & proprioceptive state, bounded to $[-5.1, 5.1]$ \\
\texttt{mode}   & 1   & commanded task mode; $1$ goal-reaching, $0$ exploration \\
\bottomrule
\end{tabular}
\end{table}

\textbf{Perception.} The \texttt{latent} block is produced by an autoencoder
trained offline and held frozen throughout policy training and evaluation. Its
identity is recorded in every checkpoint and asserted before evaluation, because
pairing a policy with a different encoder does not fail loudly---it produces
confident nonsense.

\textbf{Goal encoding.} The goal offset is expressed in the vehicle frame as the
forward and left components of the goal displacement
$d = p_{\mathrm{goal}} - p_{\mathrm{ego}}$. CARLA uses a left-handed frame with
basis vectors $f = (\cos\psi, \sin\psi)$ and $r = (-\sin\psi, \cos\psi)$, so

\begin{equation}
\label{eq:goalblock}
  \mathrm{forward} = d \cdot f,
  \qquad
  \mathrm{left} = -\,(d \cdot r).
\end{equation}

We state \eqref{eq:goalblock} explicitly because an earlier implementation
applied a textbook right-handed rotation and mirrored the lateral bearing. The
policy was thereby instructed to steer away from the goal at every turn, a
defect that manifested only as ``drives competently but never arrives.'' The
action history holds the last four applied controls ($4\!\times\!2$), supplying
the short-horizon context required for smooth control.

\textbf{Mode conditioning.} The commanded mode is an \emph{external task bit},
not a learned switch. It occupies its own observation key, and in exploration
mode the goal block is \emph{zeroed rather than omitted}, since a single actor
cannot accept two input widths. The permanently-zero channel in exploration mode
is the acknowledged cost of gating.

\subsection{Reward Functions}
\label{sec:setup:reward}

Two rewards are selected by the mode bit, drawn once per episode with
probability $0.5$.

\textbf{Goal mode: minimum time.} The goal reward is the SoRB-style
minimum-time objective,

\begin{equation}
\label{eq:sorb}
  r_t =
  \begin{cases}
    0,  & \text{on arrival (terminal)}\\
    -1, & \text{otherwise,}
  \end{cases}
\end{equation}

with an episode cap of \num{1000} steps, so the undiscounted return lies in
$[-1000, 0]$ and the optimal value function is exactly the negated time to goal.
There is no shaping, no arrival bonus, and no coefficient to balance. This is
not incidental: the study consumes the critic \emph{as a distance}, and only
under \eqref{eq:sorb} does $Q(s,a)$ read as $-(\text{steps to goal})$.

A vehicle that becomes irrecoverably wedged is billed exactly the remaining
clock, $-(1 + T_{\max} - t)$, so terminating early is worth precisely nothing.
Without this, any terminal failure would stop the per-step meter and crashing
would become optimal.

\textbf{Exploration mode: coverage.} The exploration reward pays displacement
projected onto the road direction, gated on being under target speed, against a
small per-step cost:

\begin{align}
\label{eq:explore}
  r_t = &-\,\num{4e-3} \;-\; |\Delta \mathrm{steer}| \;+\;
          g_t\,\mathbb{1}[v_t \le v^{\star}] \nonumber\\
        &-\; 3.0\,\mathbb{1}[\text{collision}]
         \;-\; 5.0\,\mathbb{1}[\text{road departure}],
\end{align}

where $g_t$ is the displacement projected onto the lane forward direction and
$v^\star$ the target speed.

\textbf{Baseline reward.} The shaped baseline receives the environment's own
per-step distance reduction plus projected progress under the same speed gate,
less terminal penalties for collision, lane departure and idling, and a per-step
steering-smoothness cost. It is the reward supplied by the upstream driving
codebase; a method that cannot beat the reward practitioners actually use has
demonstrated nothing.

\textbf{Lane-crossing convention.} Crossing a painted lane marking costs nothing
and terminates nothing, whereas leaving the drivable surface remains fatal.
Crossing markings is unavoidable when turning through junctions, and treating
every crossing as failure terminated \SI{90}{\percent} of episodes at
\numrange{44}{70} steps in an earlier configuration.

\subsection{Terminal Conditions}
\label{sec:setup:terminal}

Termination differs by mode, and the difference is material: it determines what
a success rate is comparing.

\textbf{Goal mode.} \emph{Only arrival terminates.} The environment's own
termination signal---idle, falling, distance limit, lane invasion, collision,
speed, wrong direction---is deliberately discarded, and the episode instead
truncates at the \num{1000}-step cap. Under a per-step cost, any terminal
failure would stop the meter and become an escape hatch. A wedge terminates but
is billed the full remaining clock, so it is not an escape either.

\textbf{Exploration mode.} Collision or departure from the drivable surface
terminates; the episode truncates at \num{500} steps. Off-road is retained as
fatal because, with lane crossings free, it is the only remaining term that
discourages leaving the road: the progress term in \eqref{eq:explore} pays
equally for driving down a road and for circling.

\textbf{Baseline.} The shaped baseline terminates on arrival, collision, idling,
lane invasion or off-road, and additionally honours the environment's own
termination. Because this differs from the gated arms' rule, two relaxations
allow the same frozen checkpoint to be evaluated under the gated rule: one makes
lane crossings non-terminal, and one admits arrival as the sole terminal event.
Both must suppress the environment's own termination as well as the wrapper's;
relaxing only the latter leaves the episode ending anyway, silently, because the
environment ORs its raw lane flag into its own done signal.

\subsection{Policy Learning}
\label{sec:setup:learning}

The policy is trained with Truncated Quantile Critics (TQC): five critics of
\num{25} quantiles each, dropping the two most optimistic quantiles per network,
MLP trunks of $(256, 256)$, Adam at \num{3e-4} for actor and critics and
\num{1e-5} for the temperature, $\tau = 0.005$, initial temperature $0.2$, and a
fixed target entropy of $-2.0$ for the two-dimensional action space.

\textbf{Two critics, two discounts.} Goal and exploration modes have separate
critic networks, their returns being on different scales with different
horizons. The goal critic uses $\gamma = 1.0$, which is mandatory rather than
preferred. At $\gamma = 0.99$ the effective horizon is \num{100} steps, the
value saturates near $-10$, and a \num{400}-step goal becomes indistinguishable
from a \num{1000}-step one---the two values differ by \num{0.18}---destroying
precisely the distance reading this study depends on. The undiscounted return is
safe here only because every episode terminates or truncates within the cap; it
would diverge in a genuinely non-terminating task. The exploration critic uses
$\gamma = 0.99$.

\textbf{Batch composition.} Transitions are held in a mode-partitioned replay
buffer and every batch is drawn \SI{50}{\percent}/\SI{50}{\percent} across
modes, so neither critic starves when episode lengths differ between modes
(\num{1000} steps for goal, \num{500} for exploration).

Training runs for \num{350000} environment steps in Town02 with the
goal-distance curriculum pinned to \SI{5}{\meter}--\SI{20}{\meter}. Two arms are
trained identically except for hindsight relabelling of exploration episodes
into goal-mode data; the launch scripts differ only in that flag and the output
directory, so the arms cannot drift apart in any other respect.

\subsection{Mapping Algorithm}
\label{sec:setup:mapping}

The map is a directed topological graph over states visited by the policy
itself. Construction has three stages and requires no simulator after the first.

\textbf{1) Trace collection.} The policy is driven in \emph{exploration} mode
for \num{60} episodes of \num{400} steps, recording pose, observation, and the
fields from which the goal block is computed. Exploration mode is used because
it produces coverage; the goal half drives direct lines to sampled goals and
would map only those. Candidates are thinned to one per \SI{5}{\meter} as they
are recorded, since consecutive simulator steps are $\approx\!\SI{0.28}{\meter}$
apart and the next stage is quadratic in the candidate count. This yields
\num{1510} candidates over \SI{7464}{\meter} of driving.

\textbf{2) All-pairs learned distance.} For every ordered pair $(i,j)$, the
observation at state $i$ is relabelled so that its goal block points at
candidate $j$, and the goal critic is queried:

\begin{equation}
\label{eq:dist}
  D_{ij} \;=\; -\,\frac{1}{c}\,
  Q\!\left(\,\mathrm{relabel}(o_i, p_j),\; \pi(\cdot \mid o_i)\,\right),
\end{equation}

with $c$ the per-step cost of \eqref{eq:sorb}. This requires \emph{no
simulator}: the perception latent records what the camera saw and is independent
of which goal the step is credited to. Relabelling uses the observation
wrapper's own transform, so relabelled observations cannot drift from on-policy
ones. Two constraints are load-bearing. Observations must be forced to goal mode
before any query, since the goal block is zeroed at mode $0$ and querying the
goal critic with exploration-mode inputs reads the wrong half of a gated policy.
And $c$ in \eqref{eq:dist} must be the arm's own value; a mismatch scales every
edge weight by an order of magnitude without failing.

\textbf{3) Node selection and edges.} Four placement rules
(Section~\ref{sec:setup:placement}) select nodes from the shared candidate set.
Edges are directed and weighted by $D_{ij}$. An edge is \emph{rejected} when
$D_{ij}$ exceeds the horizon $H$ on which the critic was trained
($H = \SI{30}{\meter} \approx 106$ steps): beyond $H$ the critic extrapolates,
and graph search actively selects for optimistic extrapolation error. Nothing is
symmetrised---a one-way street is genuinely asymmetric and the critic represents
that.

Critic-ensemble disagreement is deliberately \emph{not} used as an edge filter.
Measured on these agents, disagreement is a proxy for leg length rather than
local uncertainty---AUC \num{0.533}, chance, once length and $|Q|$ are regressed
out---and it carries the wrong sign, so rejecting on it would preferentially
discard the short, reliable edges.

\subsection{Localization Algorithm}
\label{sec:setup:localization}

Localization is nearest-node lookup in the plane. At episode start, the current
pose and the goal pose are each associated with the nearest graph node within
\SI{25}{\meter}. The radius is bounded by the policy's own competence rather
than tuned: its training band is \SI{5}{\meter}--\SI{20}{\meter}, so a node
within $\approx\!\SI{25}{\meter}$ is a leg the policy can be expected to drive,
and a larger radius would assert entry points it cannot reach.

Localization has three distinct failure modes, reported separately rather than
pooled into the success rate because they answer different questions: the start
pose has no node within radius, the goal pose has none, or both localize but the
graph contains no path between them. The last is a statement about graph
connectivity, not about localization.

\subsection{Node Placement Methods}
\label{sec:setup:placement}

Four rules select nodes from the \emph{same} candidate set, so placement is the
only variable. They differ in what their free parameter is, which determines how
each is swept.

\begin{itemize}
  \item \textbf{uniform}---nodes at equal arc-length spacing, taken at interval
        midpoints. It uses no learned quantity, and is the baseline a method
        must beat to justify its cost.
  \item \textbf{random}---distinct candidates drawn uniformly at random. The
        floor.
  \item \textbf{policy}---peaks of the policy's own action entropy, smoothed and
        taken greedily by height subject to a minimum separation. Entropy is
        estimated as $-\log \pi(a\mid s)$ for a single sample, the same
        estimator reported during training. The hypothesis is that the policy is
        least certain where the route is \emph{decidable}---at junctions---so
        that entropy is a free junction detector.
  \item \textbf{sgm}---Sparse Graphical Memory's two-way consistency merge, in
        which a candidate merges into an existing node when the two agree,
        within a cutoff, about their distances \emph{to} and \emph{from} every
        current node.
\end{itemize}

The first three have no intrinsic node count and are swept over the number of
nodes requested. SGM derives its own sparsity from its consistency cutoff and is
therefore swept over that cutoff, with node count as an \emph{output}. No shared
node budget is imposed, since the object of study is the Pareto frontier of node
count against navigation performance.

Two implementation notes affect reproducibility. SGM's published cutoffs are
expressed in its reference implementation's value scale and do not transfer: at
the published values nothing merges, and the algorithm silently degenerates to
retaining every candidate. Cutoffs here are derived from measured quantities on
the target critic, and a zero-merge outcome is treated as a hard error rather
than a result. Distances are additionally clipped at $H$ for the merge test,
because that test is a maximum over all nodes and would otherwise be decided by
extrapolated values.

\subsection{Evaluation Algorithm}
\label{sec:setup:eval}

\textbf{Episode set.} Evaluation uses a fixed episode set drawn from a generator
never touched by training, so every method drives identical start states and
goals. The environment is reset with each episode's own seed; omitting this
makes the fixed set decorative, since the simulator then samples freely and the
property the comparison rests on silently fails to hold.

\textbf{Navigation.} Per episode, the harness localizes start and goal
(Section~\ref{sec:setup:localization}), computes a shortest path over the
directed graph by Dijkstra with $D_{ij}$ as edge weights, and drives the node
sequence by rewriting the policy's goal block to the current waypoint. A
waypoint is advanced when the vehicle is within \SI{4}{\meter}; a leg not
completed within \num{200} steps is abandoned and the next waypoint adopted,
since a leg the policy cannot drive is an edge the graph should not have
offered. The final waypoint is always the true goal, not the last node.
Episodes for which the graph offers no path are still driven, with the true goal
alone: excluding them would score a map on the subset it happened to cover and
would flatter a sparse or poorly connected graph.

\textbf{Metrics.} Success $S_i \in \{0,1\}$ is the environment's own arrival
flag, identical to the metric used elsewhere in this codebase. Success weighted
by Path Length is

\begin{equation}
\label{eq:spl}
  \mathrm{SPL} \;=\; \frac{1}{N}\sum_{i=1}^{N}
  S_i \,\frac{\ell_i}{\max(\ell_i,\, p_i)},
\end{equation}

where $\ell_i$ is the length of CARLA's own shortest driving route between start
and goal, obtained from the map's route planner, and $p_i$ is the distance the
agent actually drove. The geodesic in \eqref{eq:spl} is a road route rather than
a Euclidean distance, since a straight line through a building is not a shortest
path and would credit every method with detours it did not avoid. Two
diagnostics are reported alongside: the number of episodes the graph could not
route, and the number of subgoal legs abandoned on the step budget.

\textbf{Undefined metrics are refused, not reported as zero.} Success is not
reported for an arm that cannot observe the goal. SPL is reported as \emph{not
applicable}, never \num{0.000}, when no road route is available, since a zero
reads as ``took an infinite detour'' rather than ``was not measured.'' Critic
calibration is reported only for arms trained on the minimum-time reward, a
shaped arm's critic predicting a shaped return so that scoring it against
elapsed steps measures nothing; calibration is computed on arriving episodes
only, a timeout having no true cost-to-go.

\subsection{Reproducibility}
\label{sec:setup:repro}

All arms share one trace, one distance matrix, one edge rule, one episode set
and one policy; placement is the only variable. Results are keyed by distance
band, since pooling an in-distribution result with an extrapolation one produces
a number describing neither. Every reported configuration is recorded in the
result file alongside its metrics---including the node count actually placed,
the consistency cutoff actually used, and the step cost---so that a row can be
checked against the configuration that produced it rather than the one that was
intended.
